% s00-front.tex —— 译本说明
\chapter*{译本说明}
\addcontentsline{toc}{chapter}{译本说明}
\markboth{译本说明}{译本说明}

\section*{底本}

本书是 David Ben-Zvi 于 2021 年春季在 UT Austin 讲授的 \textbf{M390C（Geometric Langlands）}
课程的课堂讲义中译本。讲义由听课者 Arun Debray 以 live-\TeX{} 现场记录，2021 年 6 月 3 日定稿，
共 89 页、31 讲，授课日期自 2021 年 1 月 19 日至 6 月 3 日。

底本标题为 \textit{M390C Notes: Geometric Langlands}。课程以\textbf{电磁对偶}为线索，
把几何 Langlands 纲领讲成一串拓扑场论的谱分解与对偶：从 Fourier 变换与 Pontrjagin、Cartier 对偶出发，
经拓扑量子力学、拓扑 Maxwell 理论、Tannakian 重构、Hecke 代数，到仿射 Grassmann 簇与几何 Satake 对应，
最后回到经典 Langlands 与 $L$ 函数。底本只有 1 幅嵌入图，正文以散文与行间公式为主。

译本**印刷页码与底本 PDF 页码一致**（偏移 0）。每页正文起首处以页边小字标出底本页码
（\texttt{p.\!N}），译文与底本可按页对读。

\section*{体例}

\begin{enumerate}[leftmargin=2.2em,itemsep=0.2em]
\item \textbf{编号交给 \LaTeX}。原文每讲的编号条目（Definition、Theorem、Example、Remark 等）
      共用一个计数器，自动排成「讲次.\,序号」，与底本编号一一对应，正文中不再手写编号。
\item \textbf{公式不编号}。底本的行间公式编号（如 (9.8)）在译本中一律略去；
      原文以公式号回指的句子改写为「前式」「上述映射」等描述性说法。
\item \textbf{页锚}。每页首段前有页边页码；页界落在句子中间时改用行内小号上标。
\item \textbf{译注}。楷体脚注，标「来源」等级：
      A 为原文自明或原文页码，B 为外部文献（nLab、Stacks Project、作者原文等），
      I 为本仓库既有译本已冻结的译名。每讲不超过 6 条，只写术语背景、原文省略、原文笔误三类。
\item \textbf{保留底本的口语与铅笔注}。底本是现场记录，含作者的口头插话、
      「be careful!」式提醒与 \texttt{(TODO: \ldots)} 自注。这些一律照译保留，不删不改，
      以便读者看清讲稿的原貌；其中的技术性错误另见各讲报告与
      \texttt{资料卡/04-术语文献勘误.md}。
\item \textbf{术语}。以 \texttt{工程文件/work/glossary.tsv} 为准；人名、定理名、软件名不译，
      如 Hecke、Satake、Tannakian、Cartier、Pontrjagin、Betti、de Rham、shtuka。
\end{enumerate}

\section*{致谢}

底本作者 Arun Debray 在前言中致谢 Tom Gannon、Charlie Reid、Saad Slaoui 与 Jackson Van Dyke
协助订正错漏。译本的一切错误由译者负责。
